\documentclass[10pt,a4paper]{amsart}
\usepackage[margin=19mm]{geometry}
\usepackage{amsmath,amssymb,mathtools}
\usepackage[scr=rsfs,cal=euler]{mathalfa}
\usepackage{xfrac}
\usepackage[unicode,hidelinks,bookmarks=false]{hyperref}
\newcommand{\ie}{i.e.\ }
\newcommand{\cf}{cf.\ }
\newcommand{\resp}{respectively\ }
\newcommand{\Subsection}{\subsection}
\providecommand{\nobreakdash}{\nobreak\hbox{-}}
\setlength{\emergencystretch}{2em}
\multlinegap0pt
\usepackage{bm}
\usepackage{bbm}
\usepackage{dsfont}
\usepackage{xspace}
\usepackage{cancel}
\usepackage{mathtools}
\usepackage[shortlabels]{enumitem}
\usepackage{amsthm}
\makeatletter
\@ifundefined{theorem}{\newtheorem{theorem}{Theorem}}{}
\makeatother
\newcommand{\RR}{\mathbb{R}}
\newcommand{\cH}{\mathcal{H}}
\newcommand{\cP}{\mathcal{P}}
\title[A Cram\'er-Wold theorem for mixtures]{A Cram\'er-Wold theorem for mixtures}
\author{Ricardo Fraiman, Leonardo Moreno, Thomas Ransford}
\date{Source: arXiv:2410.22038v2 (CC BY 4.0)}
\begin{document}
\maketitle

\noindent\emph{Source and licence.} A Cram\'er-Wold theorem for mixtures. Version arXiv:2410.22038v2 (CC BY 4.0), distributed under the Creative Commons Attribution 4.0 International licence, as recorded in the arXiv licence metadata for this exact version and shown on the abstract page https://arxiv.org/abs/2410.22038v2. Full rights notice: licenses/arxiv-2410.22038-CC-BY-4.0.txt.\par\smallskip

\noindent\textit{Editorial scope.} Counted unit: the Theorem whose source label is \texttt{T:mixture} in the article (source0.tex lines 209--222, source label \texttt{T:mixture}), delivered together with the article's own complete original proof of it (source0.tex lines 232--261, one \texttt{proof} environment of 30 lines). No mathematics is written, repaired or re-derived: statement and proof are the article's own text line for line. The proof is detached from the statement in the source: the article states the result at lines 209--222 and prints its proof at lines 232--261, and the proof environment's own header names the statement, so the association is the article's own. Required context retained uncounted: none. Every definition, display and label of those lines is retained unchanged. Omitted from this package and not redistributed: 1--208, 223--231, 262--811. Nothing inside a retained excerpt is omitted or altered: every retained body is delivered exactly as the article's lines are written, so no omission receipt of any kind accompanies this package. Citations: the retained excerpts carry no citation command at all, so no bibliography entry of the article is transcribed and no reference is redistributed. Reference adaptation: none was needed and none was made; every reference inside the delivered unit resolves inside it, so no unresolved reference or citation is produced. Printed slips kept literal: no printed word, symbol, hypothesis or number of the counted statement or of its proof was altered. Layout and class: the article's own class is \texttt{amsart} with packages including amssymb, amsmath, amsthm, newlfont, enumerate; the wrapper is the lane's pinned \texttt{amsart} editorial wrapper at 10pt on A4, which restates the theorem-like environments the retained text needs and adds the article's own macro definitions that the retained text needs (\texttt{\textbackslash RR}, \texttt{\textbackslash cH}, \texttt{\textbackslash cP}), with extra wrapper packages (bm, bbm, dsfont, xspace, cancel, mathtools, enumitem). The delivered numbering is the unit's own. Assets: no retained span contains a figure environment, a graphics-inclusion command, a PDF-page inclusion, a table or a TikZ picture; no image file is copied into this package, the package declares no asset dependency and no image file is redistributed with it. Licence: version arXiv:2410.22038v2 of this article is distributed under the Creative Commons Attribution 4.0 International licence, as recorded in the arXiv licence metadata for this exact version and shown on the abstract page https://arxiv.org/abs/2410.22038v2. A full rights notice is delivered at licenses/arxiv-2410.22038-CC-BY-4.0.txt. Characters: every retained excerpt is pure ASCII, so the delivered engine, XeLaTeX with its default Latin Modern text font, reports no missing character.
\begin{theorem}\label{T:mixture}
Let $\cP$ be a family of Borel probability measures on $\RR^d$,
let $\cH$ be a collection of vector subspaces of $\RR^d$, and let $m\ge1$.
Suppose that:
\begin{enumerate}[\normalfont\rm(i)]
\item\label{i:1} 
For each $H\in\cH$, the distinct measures in the set
$\{P_H:P\in\cP\}$ are linearly independent.
\item\label{i:2} 
Each partition of $\cH$ into $2m-1$ subsets contains at least one Cram\'er--Wold system for $\cP$.
\end{enumerate}
Let $P$ and $Q$ each be convex combinations of $m$ measures from $\cP$.
If $P_H=Q_H$ for all $H\in\cH$, then $P=Q$.
\end{theorem}

\begin{proof}[Proof of Theorem~\ref{T:mixture}]
We argue by contradiction. 
Suppose that $P_H=Q_H$ for all $H\in\cH$, but that $P\ne Q$. 
Then the difference $P-Q$ can be written as 
$P-Q=\sum_{j=1}^k\lambda_j P_j$, 
where $P_1,\dots,P_k\in\cP$ are distinct,  
where $\lambda_1,\dots,\lambda_k\in\RR$  are non-zero,
and where $k\le 2m$. Set
\[
\cH_j:=\{H\in\cH: (P_{1})_H= (P_{j})_H\} \quad(j=2,\dots,k).
\]
Clearly no $\cH_j$ is a Cram\'er--Wold system for $\cP$.
If we had $\cup_{j=2}^k \cH_j=\cH$, then we could construct a partition
of $\cH$ into $k-1$ sets (perhaps some of them empty) none of which is
a Cram\'er--Wold system for $\cP$, contradicting the assumption~\eqref{i:2}
on $\cH$.
We conclude that $\cup_{j=2}^k \cH_j\ne\cH$, so there exists an $H_0\in\cH$
such that $(P_{1})_{H_0}\ne (P_{j})_{H_0}~(j=2,\dots,k)$.
		
By assumption~\eqref{i:1},
the distinct measures in the set $\{(P_j)_{H_0}:j=1,\dots k\}$
are linearly independent. In particular, since $(P_{1})_{H_0}\ne (P_{j})_{H_0}~(j=2,\dots,k)$,
it follows that $(P_{1})_{H_0}$ is not in the span of $\{(P_j)_{H_0}:j=2,\dots d\}$. 
On the other hand, we have
\[
\sum_{j=1}^k\lambda_j(P_j)_{H_0}
=\Bigl(\sum_{j=1}^k\lambda_j P_j\Bigr)_{H_0}=(P-Q)_{H_0}=P_{H_0}-Q_{H_0}=0.
\]
Since $\lambda_1\ne0$, this is a contradiction.
\end{proof}

\end{document}
